\chapter{Exemple de code Tcl top-level}

\mysoftcite{Tcl} est un langage de commande et une librairie
utilisable comme interface de programmes interactifs.

\mycodefull{tcl-conex}%
{}%
{Création des structures de données pour le voisinage}%
{Les fonction \asoft{Tcl} sont associé à des fonctions \asoft{C}
  créant les objets permettant la composition des spécification de
  connexions de voisinage.}

La figure~\mygetcode{tcl-func} présente du code \asoft{Tcl} de
création et manipulation de table de transition. %
On peut y voir la liaison des règles \arule{Ink}, \arule{maj} et
\arule{anneal}. %
\mysetref{additive-rules}%
Ce sont des règles additives. %
Leurs lookup-tables peuvent être spécifiées par une table secondaire
dont le nombre d'entrées correspond au nombre de voisins. %
Les entrées de cette table secondaire codent la nouvelle valeur de la
cellule, cette valeur est soit un des états de la règle~(ici les
règles sont binaires), soit un code désignant une fonction de
génération de la nouvelle valeur. Dans le cas présent seul le code
``U'' pour \aquote{Unchanged} est utilisé. %
Pour des règles binaires nous proposons d'étudier, et avons
expérimenté la seule autre fonction possible:~``F'', pour
\aquote{Flip} bit. L'intérêt de cette approche est de permettre un
échantillonnage de l'espace des règle pour les voisinages
classiques\footnote{%
  Moore et von Neumann.}
dont la taille permet une exploration systématique. %
Le nombre de règles d'un des AC binaire pour un voisinage de Moore
reste prohibitif, même en considérant une réduction par
symétrie~\cite[]{Wolfram83} et notre plongement de règle; %
le nombre de règles additives spécifié par une table
secondaire\footnote{%
  Dont les entrées correspondent à une fonction du voisinage, la somme
  dans le cas des règles additives} %
binaire est~$2^9$ ou $2^{10}$(Selon que l'on considère le centre ou pas
dans les termes de la somme); %
le nombre de règles spécifié par un table secondaire utilisant les
deux transformations binaires possible~(U~et~F), en plus des deux
constantes d'états, est~$2^{20}$. L'idée est d'utiliser cette
technique en augmentant le nombre d'états et en modulant le choix des
fonctions de transformation figurant dans la table secondaire pour
obtenir un espace de règles dont la taille est adaptée aux ressources
disponibles fonction d'une recherche de valeurs limites pour les
mesures effectuées.

\mycodefull{tcl-data}%
{}%
{Définitions d'objets top-level}%
{Les structures de données créées correspondent aux objets décrits en
  \mygetref{space-time-law}, l'objet de type ``Simple'' est un
  séquenceur.}

\mycodetwocol{tcl-func}%
{}%
{Définitions des fonctions}%
{Les règles additives \arule{Ink}, \arule{maj} et \arule{anneal} sont
  décrites dans~\cite{Toffoli87}.}

%%
%% Local Variables:
%% mode: latex
%% mode: auto-fill
%% iso-accents-minor-mode: nil
%% TeX-master: "top"
%% TeX-command-default: "mytex"
%% Local IspellParsing: "latex-mode ~latin1"
%% Local IspellPersDict: "~/.ispell_lisp"
%% LocalWords: "lookup-tables Unchanged Flip Moore von Neumann"
%% End:
